% Fixture "twocolumn": two-column article. Anchors in the left text column must go to the left
% margin, anchors in the right column to the right margin (spec §6.4).
\documentclass[10pt,twocolumn]{article}
\usepackage[T1]{fontenc}
\usepackage[utf8]{inputenc}
\usepackage{amsmath}
\pdfinfoomitdate=1 \pdftrailerid{} \pdfsuppressptexinfo=-1
\title{Sorting Networks in Two Columns}
\author{Jonas Keller}
\date{}
\begin{document}
\maketitle
\section{Left column}
Sorting networks are comparator circuits whose structure does not depend on the input. Batcher's
odd--even merge sort uses $O(n \log^2 n)$ comparators and depth $O(\log^2 n)$. The zero--one
principle states that a network sorts all inputs if it sorts all binary inputs, which reduces
verification to $2^n$ cases. We revisit this classic argument with a short proof.

A comparator exchanges two wires when they are out of order. Composition of comparators in
layers gives the depth; the number of comparators gives the size. Lower bounds on size are
notoriously hard: only for small $n$ are optimal networks known.
\section{Right column}
In the right column we discuss the AKS network, which achieves depth $O(\log n)$ with enormous
constants. Practical implementations prefer bitonic sorters on GPUs because of their regular
memory access. The bitonic merge splits a bitonic sequence into two halves of which every
element of one is at most every element of the other.

Our main observation concerns the monotone Boolean functions computed by each output wire,
\begin{equation}
  y_k = \bigvee_{|S| = n-k+1} \bigwedge_{i\in S} x_i ,
\end{equation}
which immediately yields the zero--one principle via threshold functions.
\end{document}
